\documentclass[11pt,a4paper]{article}
\usepackage[T1]{fontenc}
\usepackage{lmodern,amsmath,amssymb,amsthm,mathtools}
\usepackage[margin=26mm]{geometry}
\usepackage{array,longtable,booktabs,xcolor,microtype}
\usepackage[colorlinks=true,linkcolor=blue!45!black,citecolor=blue!45!black,urlcolor=blue!45!black]{hyperref}
\newtheorem{theorem}{Theorem}[section]
\newtheorem{lemma}[theorem]{Lemma}
\newtheorem{proposition}[theorem]{Proposition}
\newtheorem{corollary}[theorem]{Corollary}
\theoremstyle{definition}
\newtheorem{definition}[theorem]{Definition}
\theoremstyle{remark}
\newtheorem{remark}[theorem]{Remark}
\newcommand{\Q}{\mathbb Q}
\newcommand{\Z}{\mathbb Z}
\newcommand{\Sha}{\operatorname{Sha}}
\setlength{\emergencystretch}{2em}
\setcounter{tocdepth}{1}
\title{An explicit polynomial in 29 variables over the rationals}
\author{}
\date{8 November 2026}
\begin{document}\maketitle
\begin{abstract}
We specify an integer polynomial in 29 total variables, with degree at
most $2^{174}$ and coefficient height at most $2^{\,2^{177}}$.
Under four precisely stated arithmetic inputs, its rational input
projection is exactly $\mathbb Z$. A fixed prime forces a sufficiently
large elliptic index through an existing scalar pole condition, so no
separate size certificate is needed. Nine quartic point choices and
four quadratic pole choices are eliminated by a finite-algebra
construction using one additional rational coordinate. The paper gives
the rational source, the local and real existence arguments, the
integer-polynomial formula, and reproducible degree and coefficient
bounds. The arithmetic premises are explicitly distinguished from the
algebraic and computational conclusions proved here.
\end{abstract}
\tableofcontents\newpage
\section{The construction and its scope}
For $F\in\mathbb Z[X,Y_1,\ldots,Y_m]$, its rational input projection is
\[
 \{x\in\mathbb Q:\exists\boldsymbol y\in\mathbb Q^m,
                         F(x,\boldsymbol y)=0\}.
\]
We count the input among the variables. Degree means total degree;
coefficient height means the largest absolute integer coefficient.
The polynomial constructed below has $m=28$ and satisfies
\[
                 \deg F\le2^{174},\qquad H(F)\le2^{\,2^{177}}.
\]
These are upper bounds for a specified polynomial, rather than exact
values of an expanded coefficient table. Its projection equals the
integers under the four arithmetic inputs in Section~\ref{sec:u29}.
Those inputs are results asserted in the cited upstream manuscript and
Sun's prime-pattern theorem. The implication from them is proved here;
the coefficient checks do not establish the deep arithmetic inputs.

The exact specification comprises an integer coefficient expression
graph and finite determinant, trace, resultant and product operations.
Every range is explicitly bounded. The formal roots and matrix entries
are internal to these finite operations and are not additional
existential variables.
\section{The exact numerical height constant}
\label{sec:exact-height-constant}

The contact-height input uses a fixed, explicitly specified integer.
We record its exact size for the index estimates below. Put
\[
 B_*=\prod_{p\le76176}p,\qquad
 H_{19}=2^{110}210^{1800}B_*^{19}+128.
\]
This is the constant of the source's exponent-nineteen height theorem
\cite[Theorem \nolinkurl{thm:quartic-height19}]{RepositoryV9}.

\begin{lemma}\label{lem:exact-height-constant}
The exact inequalities
\[
                 2^{2093447}<H_{19}<2^{2093448}
\]
hold. In particular one may use $h=2093448$ in every argument
requiring $H_{19}<2^h$.
\end{lemma}
\begin{proof}
This finite integer calculation has the following reproducible
certificate. Enumerate primes through $76176$ by marking, for
each unmarked $p\le\lfloor\sqrt{76176}\rfloor$, all multiples
starting at $p^2$. The remaining integers at least two are
exactly the primes: every composite has a prime divisor not
exceeding its square root, and no prime is marked.
An independent enumeration tests every integer by trial division
through its integer square root. Both lists have $7499$ entries
and agree entry by entry.

Their exact product has $109445$ binary digits. Computing the
displayed powers and product as integers gives $2093448$ binary
digits for $H_{19}$. The comparison with the next smaller power
of two is strict as well. The supplied checker performs both
prime enumerations and both integer comparisons; the complete
prime list and the unsigned binary encoding of $H_{19}$ are
included. No floating-point logarithm enters these comparisons.
\end{proof}

This lemma evaluates a constant. It does not prove the
contact-height inequality itself. The latter is still the
arithmetic input stated in Section~\ref{sec:u29}, now used with
this exact value and its sharper certified bound.


\section{An explicit polynomial for a guarded finite algebra}
\label{sec:explicit-finite-algebra}

Let $R=\mathbb Z[x_1,\ldots,x_b]$. Consider a triangular monic presentation
\[
 A=R[z_1,\ldots,z_m]/(f_1,\ldots,f_m),\qquad
 f_i\in R[z_1,\ldots,z_i],
\]
where $f_i$ is monic of degree $e_i\ge1$ in $z_i$. Its standard
monomials $\prod_i z_i^{a_i}$, $0\le a_i<e_i$, form an $R$-basis.
In particular $A$ is free of rank $n=\prod_i e_i$. For $a\in A$,
write $M_a$ for its multiplication matrix in this ordered monomial basis.
The matrices $M_a$ commute, and their entries are integer polynomials in
the base coordinates. They are completely specified by repeated reduction
using the monic equations; they need not be materialized to specify a
determinant involving them.

Fix elements $G,E_1,\ldots,E_r\in A$. We assume that in every geometric
fiber, every point at which $G$ is nonzero has local length one. This
condition holds, for example, if $G$ contains the product of the diagonal
derivatives $\partial f_i/\partial z_i$ as a factor. We seek the rational
base points admitting a rational point of $A$ with
\[
 G\ne0,\qquad E_1=\cdots=E_r=0.
\]

For $0\le\ell\le m n(n+1)/2$, define
\[
 \lambda_\ell=1+\sum_{i=1}^m\ell^i z_i,\qquad
 M=M_{G\lambda_\ell},\qquad
 H(V)=\det(I-VM)=\sum_{k=0}^n a_kV^k.
\]
Here and below $a_k\in R$ depend on the fixed integer $\ell$; always
$a_0=1$. Also put
\[
 R_i(V)=\operatorname{Tr}\bigl(M_{E_i}\operatorname{adj}(I-VM)\bigr).
\]
This is a polynomial in $R[V]$ of degree at most $n-1$.

For each $1\le d\le n$, set
\begin{align*}
 H_d(V)&=\sum_{k=0}^d a_kV^k,\\
 g&=a_d\operatorname{Res}_{V}(H_d,H'_d),\qquad L=a_d^2g,\\
 \Psi(U)&=1+\sum_{j=1}^d a_{d-j}a_d^{2j-1}g^jU^j,\\
 \widetilde R_i(U)&=(LU)^{n-1}R_i\bigl(1/(LU)\bigr),\\
 Q_{\ell,d}(\mathbf x,U)&=
 \Psi(U)^2+\sum_{k=d+1}^n a_k^2+\sum_{i=1}^r\widetilde R_i(U)^2.
\end{align*}
The resultant is the Sylvester determinant with the specified degree
bounds $d$ and $d-1$. The expression for $\widetilde R_i$ denotes
polynomial coefficient reversal, hence defines a polynomial also at
$LU=0$. Every displayed expression has integer coefficients.

\begin{theorem}\label{thm:explicit-guarded-finite-algebra}
The explicitly specified polynomial
\[
 P(\mathbf x,U)=
 \prod_{\ell=0}^{m n(n+1)/2}\ \prod_{d=1}^{n} Q_{\ell,d}(\mathbf x,U)
\]
has the following property. For every $\mathbf a\in\mathbb Q^b$,
there exists $U\in\mathbb Q$ with $P(\mathbf a,U)=0$ if and only if
there is a rational point of the fiber algebra $A_{\mathbf a}$ satisfying
$G\ne0$ and all $E_i=0$. The formula uses only one coordinate in addition
to the $b$ base coordinates; its formal algebra generators are eliminated
by the determinants and are not quantified coordinates.
\end{theorem}
\begin{proof}
Fix a rational base point. First consider one factor $Q_{\ell,d}$.
At $a_dg=0$ or $U=0$, its reciprocal residual $\Psi$ equals one.
Otherwise, with $V=1/(LU)$,
\[
 \Psi(U)=(LU)^d\frac{H_d(V)}{a_d}.
\]
Thus a rational zero of $Q_{\ell,d}$ implies $a_dgU\ne0$, all
coefficients $a_k$ with $k>d$ vanish, and $H(V)=H_d(V)=0$.
The resultant condition makes this a simple root of $H$. Since
$H(0)=1$, the root $V$ is nonzero.

The rational matrix $I-VM$ has a one-dimensional kernel. Every
multiplication matrix $M_{z_i}$ preserves it and therefore acts on
it by a rational scalar $\zeta_i$. Applying the algebra relations
to this common eigenline shows that $(\zeta_1,\ldots,\zeta_m)$ is
a rational point of the fiber. Its $G\lambda_\ell$-value is $1/V$,
so its $G$-value is nonzero.

To verify the closed conditions, write
$J=\operatorname{Tr}(\operatorname{adj}(I-VM))$. At a root of $H$,
the adjugate identity gives
$M\operatorname{adj}(I-VM)=V^{-1}\operatorname{adj}(I-VM)$.
Differentiating the determinant then yields
\[
 J=-V H'(V)\ne0.
\]
The columns of the adjugate lie in the common eigenline, so
\[
 R_i(V)=E_i(\boldsymbol\zeta)J.
\]
Since $LU\ne0$, the equation $\widetilde R_i(U)=0$ is equivalent
to $R_i(V)=0$, and hence $E_i(\boldsymbol\zeta)=0$. This proves
the reverse implication without assuming that the rest of the fiber
algebra is reduced.

Conversely, suppose a suitable rational point exists. Over an
algebraic closure, multiplication by $G\lambda_\ell$ has the
value $G(p)\lambda_\ell(p)$ at every geometric point $p$, with
multiplicity equal to its local length. If $G(p)=0$, its factor
in $\det(I-VM)$ is one, even if the local algebra is nonreduced.
If $G(p)\ne0$, its local length is one by hypothesis.

At any such point $p$, the polynomial
$G(p)(1+\sum_{i=1}^m\ell^i z_i(p))$ is nonzero as a polynomial
in $\ell$, because its constant coefficient is $G(p)\ne0$.
For distinct such points $p,q$, their difference is also nonzero:
if all coefficients vanished, first $G(p)=G(q)\ne0$, and then
all their coordinates would agree. Each of these polynomials has
degree at most $m$. With at most $n$ geometric points, there are
at most $m n(n+1)/2$ bad values of $\ell$ in total, counting zero
values and collisions. Thus one of the specified integers makes
all nonzero guarded eigenvalues distinct and nonzero.

For that $\ell$, the polynomial $H$ is squarefree of some degree
$1\le d\le n$. The given rational point provides a nonzero
rational root $V$ of $H$. Its closed residuals vanish, so the trace
identities give $R_i(V)=0$. The leading coefficient and resultant
are nonzero. Taking $U=1/(LV)$ therefore gives a zero of the
corresponding factor $Q_{\ell,d}$, and hence of $P$.
\end{proof}

The degree cases are part of the literal finite formula. They include
every exceptional drop in degree of $H$. The discriminant condition
excludes bad separating choices, and the finite bound on $\ell$ proves
that this exclusion loses no valid rational point. In particular the
construction does not leave a generic primitive-element argument or an
unperformed exceptional-stratum decomposition implicit.


\section{Odd products and a 29-variable construction}
\label{sec:u29}

Variable counts include the distinguished input.  The construction uses
Conant's odd-product height identity and direct-target comparison argument
as presented in \cite[Lemma \nolinkurl{v9:lem:product} and
Lemma \nolinkurl{lem:v9-bridge}]{RepositoryV9}, together with the explicit
finite algebra construction.  We give the rational system, prove its
input projection under precisely stated arithmetic hypotheses, and then
eliminate its finite algebraic choices.  The product identity and point
charts are proved independently of the contact-height and prime-pattern
inputs.

\subsection{Arithmetic inputs and conventions}

Put
\[
 K=\mathbb Q(\sqrt2),\quad d=3-\sqrt2,\quad D=d^2=11-6\sqrt2,
 \qquad E:y^2=x^3-Dx,
\]
and fix
\[
 R_0=(1+\sqrt2,2\sqrt2),\qquad
 P=2R_0=\left(\frac{57}{8}-\frac72\sqrt2,
          -\frac{165}{8}+\frac{417}{32}\sqrt2\right).
\]
Write $\sigma$ for conjugation, $T_0=(0,0)$, $f(Q)=1/x(Q)$,
and $z(Q)=-x(Q)/y(Q)$.  Heights are absolute logarithmic Weil
heights, with normalized local degree weights.  Valuations at
unramified primes satisfy $v(p)=1$.

We use the following arithmetic hypotheses, with the named source
locations in \cite{RepositoryV9}.
\begin{enumerate}
\item (Elliptic group and height.)\\
      Source: \nolinkurl{subsec:v5-point-source}, \nolinkurl{eq:o8-indexheight}.\par $2E(K)=\mathbb ZP$, and
      $h(x([n]P))\ge n^2-6$ for every nonzero integer $n$.
\item (Contact height: Theorem \nolinkurl{thm:quartic-height19}.)
      For $t=u/v$ in lowest terms, $v>0$, and
      $t\notin\mathcal D=\{0,1,2401,-343/32\}$, define $M_5(t)$
      to be the product of the primes $p>76176$ for which at least one of
\[
 v_p(v),\quad v_p(u),\quad v_p(u-v),\quad
 v_p(u-2401v),\quad v_p(32u+343v)
\]
      is positive and odd.  These are the normalized contacts with
      $\infty,0,1,2401,-343/32$.  Put
\[
 B_* =\prod_{p\le76176}p,\qquad
 H=H_{19}=2^{110}210^{1800}B_*^{19}+128.
\]
      The inherited assertion is $h(t)\le H M_5(t)^{19}$.
      Exact integer evaluation gives $H<2^{2093448}$.
      We use this exponent without weakening it, so $\eta=19$.
\item (Pole obstruction and realization.)\\
      Source: \nolinkurl{hyp:o8-pole} and
      \nolinkurl{lem:original-seven-simultaneous-realization}, and
      \nolinkurl{v5:preparation-zero-bridge}.
      The pole equations, with $B=17b$,
\[
 U=b\frac{1-bt^2}{1+bt^2},\quad h_0=e(B-U^2),\quad
 p^2=B(1-e)(1+3e),\quad
 Bq^2=(s^2-Bh_0)(s^2+3Bh_0),
\]
      exclude negative odd $v_\ell(b)$ for $\ell\notin\{2,3,17\}$.
      They are soluble with $t=0$ whenever $b$ has square reduced
      denominator, is positive and less than $17$, is square in
      $\mathbb Q_2,\mathbb Q_3,\mathbb Q_{17}$, and satisfies
      $v_2(b)\ge3$, $v_3(b)\ge1$, $v_{17}(b)\ge2$.
\item (Simultaneous prime patterns: Theorem \nolinkurl{thm:v5-simultaneous-primes},
      Lemma \nolinkurl{lem:v5-two-reductions}, and the application of
      \cite[Theorem 4.1]{Sun2026Rational86}.)
      The simultaneous Gaussian prime-pattern input supplies the
      prescribed prime values $u+c_jw^2$, arbitrary fixed divisibility
      of $w$, and arbitrarily large $u/w^2$.  Their two conjugate
      reduction orders satisfy
\[
 v_2(d_-)=6,\quad v_2(d_+)=2,\quad
 \gcd(d_-(p),d_+(p'))=4
\]
      for every pair of selected primes.  An arbitrary fixed odd
      avoidance modulus can be imposed.  For each selected Gaussian
      prime $\pi$, the source convention is
      $\pi\equiv1+3i\pmod7$, making both $N(\pi-1)$ and
      $N(\pi+1)$ units at seven.  At each other odd prime $\ell$
      in the fixed avoidance modulus, the condition
      $\pi\equiv i\pmod\ell$ makes both orders congruent to two
      modulo $\ell$.  Thus the orders avoid the entire odd part
      of the fixed modulus.  The selected rational primes are $1$
      modulo $8\cdot3\cdot17$.
\end{enumerate}
These hypotheses include the simultaneous application of Sun's stated
prime-pattern theorem.  They are assumptions of the theorem below,
not conclusions of its implementation audit.
The corrected arithmetic proof of the pole realization is the one
in Version~11 or later of the source manuscript \cite{RepositoryV11}.
The finite-avoidance and zero-preparation consequences are identified
precisely in the accompanying arithmetic dependency paper.

The assigned comparison orders are $\mu=10$ and $\kappa=39$.
The logarithm truncation parameters are $\rho_0=12$ and $\rho_1=40$,
and the scalar ratio order is $m=20$. The fixed-prime argument in
Section~\ref{sec:fixed-prime-contact} supplies the needed index size.

\subsection{An elementary Lucas growth estimate}

The following estimate is independent of the arithmetic hypotheses.
It supplies the growth inequality used in the integer-extraction argument.

\begin{lemma}\label{lem:lucas-growth-elementary}
Let $a,t,b\in\mathbb Z$, with $|a|>7$ and $b\ne0$, satisfy
\[
 t^2-a(a-2)tb+(a-2)^2b^2=1.
\]
Then $\log|t|>|a|-3$. In particular,
$\log\max(|a|,|t|,|b|)>|a|-3$.
\end{lemma}
\begin{proof}
Put $A=|a|$, $c=(a-2)b$ and $c'=\operatorname{sgn}(a)c$.
Then
\[
                       t^2-Atc'+(c')^2=1.
\]
We have $|c'|\ge|a-2|\ge6$, so $t\ne0$. The two terms $t,c'$
have the same sign: opposite signs would make the left-hand side
at least $t^2+(c')^2>1$. Thus $T=|t|$ and $C=|c'|$ are positive
integer solutions of $T^2-ATC+C^2=1$.

Define $U_0=0$, $U_1=1$ and $U_{j+1}=AU_j-U_{j-1}$.
Every positive solution consists of adjacent terms of this sequence.
Indeed, order its coordinates as $M\ge m>0$. The integer
\[
                     Am-M=\frac{m^2-1}{M}
\]
lies in $[0,m)$, and replacing $(M,m)$ by $(m,Am-M)$ preserves
the equation. Repeated descent ends at $(1,0)$. Reversing the
steps gives the unordered pairs $\{U_{j+1},U_j\}$. Consequently
$C=U_j$ for some $j\ge1$ and $T\ge U_{j-1}$.

Put $d=|a-2|$. Since $d\mid C$, reduction of the recurrence
modulo $d$ gives $d\mid j$: if $a>0$, then $A\equiv2\pmod d$
and $U_j\equiv j\pmod d$; if $a<0$, then $A\equiv-2\pmod d$
and $U_j\equiv(-1)^{j-1}j\pmod d$. Hence $j\ge d\ge A-2$.
The recurrence also gives $U_r>(A-1)^{r-1}$ for $r\ge2$.
This follows by induction from $U_2=A>A-1$ and
$U_{r+1}>(A-1)U_r$, since $U_r>U_{r-1}$.
As $j\ge6$, we obtain
\[
 \log T>(j-2)\log(A-1)\ge(A-4)\log7>A-3
 \qquad(A\ge8).
\]
The last inequality holds at $A=8$ and its difference increases
with $A$. This proves the assertion.
\end{proof}


\subsection{An exact height identity}

\begin{lemma}\label{lem:u29-product}
For any nonzero odd integers $n_1,\ldots,n_r$,
\[
 h\left(\prod_{i=1}^r x([n_i]P)\right)
       =\sum_{i=1}^r h(x([n_i]P)).
\]
This identity requires only that the displayed $P=2R_0$ have infinite
order; it does not require the contact-height or prime-pattern hypotheses.
\end{lemma}
\begin{proof}
First consider a finite place $v$ with $v(D)=0$.  The formal subgroup
$E_1(K_v)$ consists of the identity and points with $v(x)<0$.
It is an open subgroup, so $E(K_v)/E_1(K_v)$ is finite.  These formal-group facts hold for the displayed integral Weierstrass
equation, including the place above two \cite[Chapter~VII]{Silverman2009}.  The image of $T_0$ has order two, since
$T_0\notin E_1(K_v)$, and
\[
                    x(Q+T_0)=-D/x(Q).
\]
If an odd multiple of $P$ belongs to $E_1(K_v)$, the image of $P$
in the quotient has odd order.  No multiple of that image is the
nonzero element represented by $T_0$.  Hence no odd multiple of $P$
has positive $x$-valuation.  If no odd multiple belongs to $E_1(K_v)$,
all odd multiples have nonnegative $x$-valuation.  Thus the numbers
$v(x([n_i]P))$ have a common weak sign.

The only remaining finite place is the prime $(3-\sqrt2)$ above seven.
Here $v(D)=2$, and the discriminant has valuation six, so the displayed
integral model is minimal.  Its reduction is $y^2=x^3$.
The displayed $P$ reduces to $(1,6)$, a nonsingular point.  The subgroup
of points with nonsingular reduction therefore contains every multiple
of $P$.  A point with positive $x$-valuation would have positive
$y$-valuation and reduce to the singular point $(0,0)$.  Consequently
every multiple of $P$ has $x$-valuation at most zero at this place.

At either real embedding every $[n_i]P$ is a nonzero double.  The
duplication identity
\[
 x(2Q)=\left(\frac{x(Q)^2+D}{2y(Q)}\right)^2
\]
makes its abscissa positive.  The curve equation and its nonzero
ordinate imply $x([n_i]P)>d>1$ in each embedding.  Each real local
height therefore adds under multiplication.  At a finite place the
common weak sign similarly gives
\[
 \max\{0,-\textstyle\sum_i v(x([n_i]P))\}
       =\sum_i\max\{0,-v(x([n_i]P))\}.
\]
Summing with the normalized local weights proves the identity.
\end{proof}

\begin{corollary}\label{cor:u29-scales}
Let $Q_1,\ldots,Q_4$ be odd multiples of $P$ and write
\[
 C=\prod_{i=1}^4x(Q_i)=C_0+C_1\sqrt2.
\]
The three numbers $A_0=1$, $A_1=C_0^2$, and
$A_2=(1+C_1^2)^2$ are nonzero positive rational squares.
If the three contact detectors at $A_j\lambda$ imply
$M_0(A_j\lambda)^k\mid N\ne0$, then
\[
                         h(x(Q_i))<7H|N|^{\eta/k}
                         \qquad(1\le i\le4).
\]
Here $\Theta=\{0,-1/2,-1,1/2399,-32/407\}$, and $M_0(t)$
is the product of the primes $p>76176$ for which $v_p(t-\theta)$
is positive and odd for at least one $\theta\in\Theta$.
\end{corollary}
\begin{proof}
The product is positive in both embeddings, so $C_0>0$.
The fractional linear map $t\mapsto2+1/t$ identifies the five
contacts with the original five, outside the exceptional primes.
Consequently $h(t)\le(H+\log3)M_0(t)^{\eta}$.  Put
$T=(H+\log3)|N|^{\eta/k}$.  The unscaled detector and each scaled
detector give $h(A_j)\le2T$.  Hence $h(C_0)\le T$.
For $C_1=u/v$ in lowest terms with $v>0$,
$h(1+C_1^2)=\log(u^2+v^2)\ge2h(C_1)$, so $h(C_1)\le T/2$.
The elementary height inequality in the fixed basis gives
\[
 h(C)\le\tfrac32T+\tfrac32\log2<7H|N|^{\eta/k}.
\]
Lemma~\ref{lem:u29-product} bounds every summand by $h(C)$.
\end{proof}

\subsection{Point charts with no free parity witnesses}

\begin{lemma}\label{lem:u29-chart}
Suppose $R=(a,b)\in E(K)$ is not a two-torsion point.  At least one
of $R,R+T_0$ has nonzero rational component of its ordinate.
Thus every nonzero double is represented by $2R$ in the chart below,
and every odd multiple $[n]P$ with $n\ne1$ is represented by
$P+4R$ in the same chart.
\end{lemma}
\begin{proof}
If both ordinates had zero rational component, then
$b\in\sqrt2\mathbb Q^\times$ and
$D/a^2=b(R+T_0)/b(R)\in\mathbb Q$.  Therefore
$a/d\in\mathbb Q\cup\sqrt2\mathbb Q$.  Since
$b^2=d^3(a/d)((a/d)^2-1)\in\mathbb Q^\times$, this would force
$d^3\in\mathbb Q\cup\sqrt2\mathbb Q$, contrary to
$d^3=45-29\sqrt2$.  Translation by $T_0$ preserves both $2R$ and
$4R$.

If $n$ is odd and $n\ne1$, take $R=[(n-1)/2]R_0$, so
$P+4R=[n]P$.  Its $4R$ is nonzero.  The affine addition denominator
$x(4R)-x(P)$ cannot vanish: otherwise $4R=\pm P$, whereas
$4E(K)=2\mathbb ZP$ contains only even multiples.  The same
translation by $T_0$ ensures the chosen chart is defined.
\end{proof}

For each of the nine names $A,S,T,Z,H_0,U_A,U_S,U_T,U_Z$, retain a
rational pair $(a_0,a_1)$.  Write $a=a_0+a_1\sqrt2$ and
$a^3-Da=A'+B'\sqrt2$.  Introduce a finite algebra root $Y$ with
\[
 Y^4-4A'Y^2+8(B')^2=0,\qquad
 b=Y/2+(B'/Y)\sqrt2.
\]
Require $Y\ne0$, $(A')^2-2(B')^2\ne0$, and
$4Y^3-8A'Y\ne0$.  Then $R=(a,b)\in E(K)$.
For the four root names set the output point to $P+4R$; for the other
five set it to $2R$.  All arithmetic is by the usual affine addition
and duplication formulas with their denominators guarded.  The
derivative condition is automatic at rational points of this chart
but is retained to ensure geometric simplicity in the finite algebra.
Indeed, since $Y\ne0$, a vanishing derivative would give $Y^2=2A'$;
the quartic would then give $(A')^2=2(B')^2$, forcing $A'=B'=0$
over $\mathbb Q$, contrary to the norm guard.
It introduces no extra coordinate.

In particular, put $C_S=S-P=4R_S$.  This point, rather than $S$,
will be used in the scalar pole certificate.  Write
\[
 (A,S,T,Z)=([n]P,[a]P,[t]P,[b]P).
\]
The four labels are nonzero odd integers; the labels are consequences
of the group hypothesis, not additional quantified variables.

\subsection{The rational system}

Retain the input $X$, the eighteen point-chart coordinates, one
rational $\lambda$, two parameters for each of three comparison poles,
and two parameters for a scalar pole.  There are
\[
                            1+18+1+6+2=28
\]
retained rational coordinates.  Set
\[
                 \beta=216X^4+450459.
\]
The following elementary input map eliminates every scalar correction
choice.

\begin{lemma}\label{lem:u29-input-map}
For rational $X$, the number $\beta=216X^4+450459$ is integral
if and only if $X$ is integral. It is greater than $2401$.
For integral $X$, its contact
$r=\beta(\beta-1)(\beta-2401)(32\beta+343)(2\beta-1)$
is a square in $\mathbb Q_2,\mathbb Q_3,\mathbb Q_{17}$,
with $v_2(r)=v_3(r)=2$ and $v_{17}(r)\in\{0,2\}$.
\end{lemma}
\begin{proof}
This is proved, with its complete residue table and exceptional
valuation calculation, in Lemmas~\ref{sn:integral-input}
and~\ref{sn:local-squares}.
\end{proof}

Put $Q_0=P$, $(Q_1,Q_2,Q_3,Q_4)=(A,S,T,Z)$,
$U_0=H_0$, and $V_i=U_i-Q_i$ for $0\le i\le4$.
For $k\in\{\rho_0,\rho_1\}$ define
\[
 \ell_k(Q)=z(Q)\sum_{j=0}^{k/4-1}c_j(Df(Q)^2)^j,
 \quad c_0=1,\quad c_1=-2/5,\quad
 c_j=c_{j-1}\frac{4j-5}{4j+1}.
\]
Put $(k_1,k_2,k_3,k_4)=(\mu,\kappa,\kappa,\kappa)$ and
$(r_1,r_2,r_3,r_4)=(\rho_0,\rho_1,\rho_1,\rho_1)$.  Set
\[
 L_i=\frac{\ell_{r_i}(U_i)}{\ell_{\rho_1}(H_0)},\qquad
 L'_i=\sigma\left(\frac{\ell_{r_i}(V_i)}{\ell_{\rho_1}(V_0)}\right).
\]
The sixteen field-valued tests, with their orders, are
\[
 \begin{array}{c|c}
 f(U_0),\ldots,f(U_4),\sigma f(V_0),\ldots,\sigma f(V_4)&1\text{ each}\\
 L_i-L'_i&k_i\\
 L_1-\beta&\mu\\
 L_3^2-L_2(L_2-2)L_3L_4+(L_2-2)^2L_4^2-1&\kappa.
 \end{array}
\]
Denote these tests by $D_j$ and their orders by $m_j$.
Use precisely the three scales of Corollary~\ref{cor:u29-scales}.
For each such $A_r$ put
\[
 t_r=A_r\lambda,\quad \delta_r=t_r-1,\quad
 F_r=\prod_{\theta\in\Theta}(t_r-\theta),\quad
 R_r=\frac{F_r}{A_r^5\delta_r^6},
\]
\[
 W_r=\prod_{j=1}^{16}
 \left(1+R_r\left(\frac{D_j}{\delta_r^{m_j+1}R_r^{m_j}}\right)^2\right),
 \qquad b_r=\frac{R_r}{\operatorname{Tr}_{K/\mathbb Q}(W_r^{-1})^2}.
\]
All $F_r$, $\delta_r$, norms of $W_r$, and displayed denominators
are required to be nonzero.  In the coefficient DAG, if
$W_r=w_0+w_1\sqrt2$, the last expression is exactly
\[
 b_r=\frac{F_r(w_0^2-2w_1^2)^2}
              {4A_r^5\delta_r^6w_0^2}.
\]

For each nonzero pole input $b$ use two retained parameters $\tau,s$
and a fixed finite choice $\varepsilon\in\{0,1\}$.  Set
\[
 B=17b,\quad e=\frac{\tau^2+(4\varepsilon-1)B}{\tau^2+3B},
 \quad h_0=e(B-b^2),\quad
 q^2=\frac{(s^2-Bh_0)(s^2+3Bh_0)}{B}.
\]
Require $\tau^2+3B\ne0$.  The first original conic is recovered by
$p=-4B\tau(1-\varepsilon)/(\tau^2+3B)$.
Conversely every point of that conic is covered: if $e\ne1$ take
$\tau=p/(e-1)$ and $\varepsilon=0$; if $e=1$ take
$\varepsilon=1,\tau=0$.  This includes $e=-1/3,p=0$.
There is no nonzero condition on $\tau,p$, or $q$.

For the fourth, scalar pole put
\[
 r_\beta=\beta(\beta-1)(\beta-2401)(32\beta+343)(2\beta-1).
\]
\[
 \delta_\beta=1+(17\ell^k r_\beta)^2,\quad
 R_\beta=\frac{102^2\ell^{2k+1}r_\beta}{\delta_\beta^2}.
\]
Here $\ell=76249$ and $k=137715$; the same constants are fixed in
every occurrence of the scalar contact.
Require $R_\beta\ne0$.  Write the rational components of
$f(H_0),f(C_S),z(C_S)/z(H_0)$ as $g_0,\ldots,g_5$, in that order.
For $N_R(u,v)=u^2+Rv^2$, set
\[
 w_0=N_R(1,g_0/R),\quad
 w_1=N_R(g_1/R,g_2/R),\quad w_2=g_3/R,
 \quad w_3=N_R(g_4/R^m,g_5/R^m),
\]
\[
              W=N_R\bigl(N_R(w_0,w_1),N_R(w_2,w_3)\bigr),
              \qquad R=R_\beta.
\]
Require $W\ne0$, and give $R_\beta W^2$ the same pole certificate.
There is no scalar correction parameter.  The only original finite
choices are the four conic bits, hence $2^4$ source branches.
For rational input $R_\beta>0$, the constant-left norm tree in fact
has $W\ge1$, so this nonzero condition is automatic; it is harmless
and convenient to retain the same denominator-guard convention.

\subsection{Local implications and soundness}

We include the local argument to make the effect of direct targets
explicit.  At a prime $p>76176$ where $v_p(R)$ is odd, consider
$W=\prod_j(1+R(D_j/R^{m_j})^2)$, $T=\operatorname{Tr}(W^{-1})$,
and a pole certificate for $R/T^2$.  Each factor has valuation
\[
 \min\{0,v_p(R)+2(v_{\mathfrak p}(D_j)-m_jv_p(R))\},
\]
since its two terms have valuations of different parity.  If
$\lambda=\max_{\mathfrak p\mid p}v_{\mathfrak p}(W)\le0$, then
$v_p(T)\ge-\lambda$.  The pole implication gives
\[
                 0<v_p(R/T^2)\le v_p(R)+2\lambda.
\]
Thus $v_p(R)>0$ and $\lambda>-v_p(R)/2$.  At a place attaining
the maximum, every factor valuation is at least $\lambda$, whence
$v_{\mathfrak p}(D_j)>(m_j-3/4)v_p(R)\ge m_j-3/4$ and therefore
$v_{\mathfrak p}(D_j)\ge m_j$.
For the normalized detectors, $\delta_r$ is a unit at each consumed
contact prime, so this argument applies to the original $D_j$.

Write $H_0=[h]P$ and $U_i=[u_i]P$.  The first ten tests place the
two families in the corresponding formal groups.  Formal logarithms
have exact ratios $u_i/h$ and $(u_i-n_i)/(h-1)$.
For the two truncation lengths $J=3,10$, the relative errors have
valuation at least $4J$. Here is a direct denominator check. The exact
formula $c_j=2\prod_{i=1}^j(4i-5)/(4i+1)$ gives
$v_p(c_j)\ge-\lfloor\log_p(4j+1)\rfloor\ge-j$ for every odd $p$:
the two progressions have counts differing by at most one at each
prime power. Terms with $j\ge\lceil4J/3\rceil$ therefore have
relative valuation at least $3j\ge4J$. For the remaining omitted
terms $J\le j<\lceil4J/3\rceil$, every denominator factor is less
than $p>76176$, so their valuations are at least $4j\ge4J$.
Thus both ratios have the required precision $k_i$.  At least one of $h,h-1$ is a $p$-unit.
The associated ratio is integral; the cross-test of order $k_i$
makes both ratios integral.  Relative precision is therefore absolute
precision.  Multiplication by $h$ and $h-1$ and subtraction gives
\[
                  L_i\equiv n_i\pmod{p^{k_i}}.
\]
This argument also holds at an inert prime.  Integer polynomial
congruences then give the corresponding congruence for the direct
Lucas residual.

The scalar certificate gives
\[
                 M_5(\beta)^m\mid(a-1),\qquad
                 h(\beta)\le H|a-1|^{\eta/m}.
\]
Indeed negative valuations of $\beta=216X^4+450459$
are multiples of four outside the fixed primes two and three.
At every consumed finite odd contact, $R_\beta$ has positive odd
valuation $\nu$; its normalizer and fixed multiplier are units. The prime
$\ell$ is never a contact counted by $M_5(\beta)$, as proved below.
For $N_R$ there is no cancellation between the two terms.  If a leaf
at depth $j$ contributes to the root, its valuation has the form
$2^jv(\text{leaf})+e\nu$, $0\le e<2^j$.
The pole implication forces the root valuation to exceed $-\nu/2$;
hence every normalized leaf has valuation greater than $-\nu$.
An original test assigned order $m$ consequently has valuation
at least $(m-1)\nu+1\ge m$.  The order-one tests put $H_0,C_S$
in the formal group; the order-$m$ ratio implies
$v_p((a-1)/h)\ge m$, and thus $p^m\mid a-1$.
The contact-height hypothesis proves the stated height bound.

Set $c=|a-1|$, $B=\max(|n|,|a|,|t|,|b|)$ and
$\Delta=t^2-a(a-2)tb+(a-2)^2b^2-1$.
If $\Delta\ne0$, the local implications and the three product-height
scales give
\[
                    B^2-6<7H(5B^4)^{\eta/\kappa}.
\]
The index size forced in Section~\ref{sec:fixed-prime-contact}
contradicts this inequality.  Hence $\Delta=0$
and $\log B>|a|-3\ge c-4$.
If $\beta\ne n$, let $N$ be the absolute value of the numerator of
$\beta-n$ in lowest terms.  It satisfies
$N\le2e^{h(\beta)}B$, and each contact radical has its $\mu$th power
dividing $N$.  Therefore
\[
                   B^2-6<7H(2e^{h(\beta)}B)^{\eta/\mu}.
\]
The scalar estimate and the inequalities in that section again contradict
Lucas growth.  Thus $\beta=n\in\mathbb Z$, and
Lemma~\ref{lem:u29-input-map} implies $X\in\mathbb Z$.

\section{A fixed-prime contact and a construction with 29 variables}
\label{sec:fixed-prime-contact}

The size of the scalar elliptic index can be forced by the scalar pole
condition itself. This removes the separate size certificate. We give
all constants and the local argument explicitly.

Throughout this section let
\[
 K=\mathbb Q(\sqrt2),\qquad D=11-6\sqrt2,\qquad
 E:y^2=x^3-Dx,
\]
and let $P$ be the normalized elliptic source point. Put
\[
 H=2^{110}210^{1800}\left(\prod_{p\le76176}p\right)^{19}+128,
 \qquad H<2^h,\qquad h=2093448.
\]
The four arithmetic inputs are the cited elliptic-group and height
results, the seven-equation pole result, and Sun's simultaneous
prime-pattern theorem. The use of those inputs, including arbitrary
fixed finite avoidance and zero preparation, is established in the
arithmetic dependency argument. No additional global existence
assertion is introduced below.

\subsection{An explicit prime with no contact root}

For the rational input $X$ define
\[
 \beta=216X^4+450459,\qquad
 r=\beta(\beta-1)(\beta-2401)(32\beta+343)(2\beta-1).
\]
In particular $\beta>2401$ and $r>0$ for every rational $X$.
Fix
\[
                \ell=76249,\qquad k=137715.
\]
For $\theta=0,1,2401,-343/32,1/2$, respectively, the five elements
$(\theta-450459)/216$ of $\mathbb F_\ell$ are
\[
                     32862,\ 32509,\ 24048,\ 904,\ 70810.
\]
Their powers with exponent $(\ell-1)/4$ are respectively
\[
                     76248,\ 76248,\ 76248,\ 76248,\ 40177.
\]
None is $1$ or $0$. The integer $\ell$ is prime: division by every
integer from $2$ through $276=\lfloor\sqrt\ell\rfloor$ gives a
nonzero remainder. These finite calculations, including the complete
trial-division test, are reproduced by the supplied independent integer
checker. They imply that none of the five equations $\beta=\theta$
has a solution in $\mathbb F_\ell$. Notice that the assertion needed
here is failure to be a fourth power, not failure to be a square.

\begin{lemma}
\label{lem:fixed-prime-valuation}
For every $X\in\mathbb Q$,
\[
 v_\ell(r)=
 \begin{cases}
 0,&v_\ell(X)\ge0,\\
 20v_\ell(X),&v_\ell(X)<0.
 \end{cases}
\]
Define
\begin{equation}
 R=\frac{102^2\ell^{2k+1}r}{\bigl(1+(17\ell^k r)^2\bigr)^2}.
 \label{eq:fixed-prime-R}
\end{equation}
Then $\nu=v_\ell(R)$ is odd and satisfies
\[
                         \nu\ge k+16=137731.
\]
\end{lemma}
\begin{proof}
If $X$ is $\ell$-integral, the preceding five fourth-power tests show
that every factor of $r$ is an $\ell$-adic unit. If $v_\ell(X)<0$,
the term $216X^4$ strictly dominates the constant term in every factor.
The coefficients $216,2,32$ are $\ell$-adic units, so each factor has
valuation $4v_\ell(X)$.

Write $t=v_\ell(r)\in20\mathbb Z_{\le0}$. Because
$k\equiv15\pmod{20}$, the number $k+t$ is nonzero. Consequently the
two summands of $1+(17\ell^k r)^2$ have different valuations, and
\[
 \nu=2k+1+t-4\min(0,k+t).
\]
This is odd. Write $k=20q+15$ and $t=-20j$, $j\ge0$. For $j\le q$
the expression is $2k+1-20j$, decreasing to $k+16$. For $j\ge q+1$
it is $1-2k+60j$, increasing from $k+16$. This proves the stated
minimum for every rational input, including inputs with an
$\ell$-adic denominator.
\end{proof}

\subsection{The scalar norm forces the necessary index size}

Set $C_S=S-P=[s]P$ and $H_0=[u]P$. Both points are nonzero in the
source, so $s,u$ are nonzero integers. As usual write $f=1/x$ and
$z=-x/y$. Let $g_0,\ldots,g_5$ be, in order, the rational components
of $f(H_0)$, $f(C_S)$ and $z(C_S)/z(H_0)$ in the basis
$1,\sqrt2$. For $N_R(v,w)=v^2+Rw^2$ set
\[
 \begin{aligned}
 w_0&=N_R(1,g_0/R),&w_1&=N_R(g_1/R,g_2/R),\\
 w_2&=g_3/R,&w_3&=N_R(g_4/R^{20},g_5/R^{20}),
 \end{aligned}
\]
\[
                  W=N_R\bigl(N_R(w_0,w_1),N_R(w_2,w_3)\bigr).
\]
The scalar pole input is $RW^2$, with preparation parameter zero.

\begin{lemma}
\label{lem:fixed-prime-index}
The scalar pole condition implies
\[
                v_\ell(s)\ge2616890,
                \qquad |s|>(64H)^{20}.
\]
\end{lemma}
\begin{proof}
Because $\nu=v_\ell(R)$ is odd, the terms $v^2$ and $Rw^2$ have
different valuations whenever both are nonzero. Thus a norm node has
valuation equal to the smaller of the two term valuations. As $\ell$
is larger than $76176$, the pole obstruction applied to $RW^2$ gives
\[
                         v_\ell(W)>-\nu/2.
\]
Every leaf contribution at depth $j$ has valuation
$2^jv_\ell(\text{leaf})+e\nu$ with $0\le e<2^j$.
It follows that every normalized leaf has valuation greater than
$-\nu$. In particular
\[
                 v_\ell(g_4),\ v_\ell(g_5)>19\nu.
\]
The order-one leaves similarly have positive valuation and put both
$H_0$ and $C_S$ in the formal group at every place above $\ell$.
The formal logarithm divided by $z$ is a unit at such a point. Its
exact scalar ratio is $s/u$, so
\[
 v_{\mathfrak l}\bigl(z(C_S)/z(H_0)\bigr)
                  =v_\ell(s/u).
\]
The components of this ratio both have valuation greater than
$19\nu$, hence $v_\ell(s/u)\ge19\nu+1$. Since $u\in\mathbb Z$,
\[
 v_\ell(s)\ge19\nu+1\ge19\cdot137731+1=2616890.
\]
Finally $\ell>2^{16}$ and
\[
 16\cdot2616890=41870240
       >41869080=20(h+6).
\]
Therefore $|s|\ge\ell^{2616890}>2^{20(h+6)}>(64H)^{20}$.
\end{proof}

For clarity, the numerical implication from this index size to
soundness can be checked directly. Use assigned comparison orders
$\mu=10$ and $\kappa=39$, scalar order $m=20$, and formal-logarithm
precisions $12$ and $40$. Put $c=|s|=|a-1|$ and
$B=\max(|n|,|a|,|t|,|b|)$. The local detector and height arguments give
\begin{align*}
 h(\beta)&\le Hc^{19/20},\\
 \Delta\ne0&\ \Longrightarrow\ B^2-6<7H(5B^4)^{19/39},\\
 \beta\ne n&\ \Longrightarrow\ B^2-6<7H(2e^{h(\beta)}B)^{19/10},
\end{align*}
where $\Delta=t^2-a(a-2)tb+(a-2)^2b^2-1$. The first inequality and
$c>(64H)^{20}$ imply $h(\beta)<c/64$.
If $\Delta\ne0$, then $5^{19/39}<\sqrt5<16/7$ and $B>\sqrt{12}$
give $B^{2/39}<32H$. But
\[
 B\ge |a|\ge c-1>c/2>(64H)^{20}/2>(32H)^{39/2},
\]
a contradiction. Hence $\Delta=0$, and the proved elementary Lucas
growth lemma gives $\log B>c-4$.
If $\beta\ne n$, the last displayed implication gives
\[
 \log B<10\log(64H)+19h(\beta)<29c/64<c-4,
\]
again a contradiction. In the second inequality we used
$\log(64H)<c/64$, which follows from $H\ge1$ and
$c>(64H)^{20}$. Thus $\beta=n\in\mathbb Z$ and the input-map
integrality lemma gives $X\in\mathbb Z$.

The scalar height argument used here still applies at every prime
counted by $M_5(\beta)$. Indeed, $\ell$ is never such a prime: at
$\ell$ the input either has no contact or its denominator valuation
is a multiple of four. At any other consumed contact prime, the extra
factor $\ell^{2k+1}$ is a unit and the new normalizer is a unit. The
valuation of $R$ is therefore the same positive odd valuation as for
$r$. The norm-tree argument gives $p^{20}\mid s$, as required.

\subsection{Realization of the scalar pole for every integer input}

\begin{lemma}
\label{lem:fixed-prime-completeness}
For every integer $X$, the scalar normalization in
\eqref{eq:fixed-prime-R} permits the simultaneous source construction
with zero preparation. It adds no retained variable, finite generator,
or independent arithmetic hypothesis.
\end{lemma}
\begin{proof}
For integer $X$ the input-map congruences give a positive integer $r$
which is a square in $\mathbb Q_2$, $\mathbb Q_3$ and
$\mathbb Q_{17}$, with
\[
                     v_2(r)=v_3(r)=2,\qquad v_{17}(r)\in\{0,2\}.
\]
Since $\ell\equiv1\pmod{24}$ and $\ell\equiv4\pmod{17}$, the
integer $\ell$ is itself a square in each of these three local fields.
The normalizer $\delta=1+(17\ell^k r)^2$ is an integer and a unit at
$2,3,17$. Moreover $\gcd(r,\delta)=1$ and $\delta\equiv1\pmod\ell$,
so its denominator has no common prime factor with the numerator of $R$.
Thus $R$ has square reduced denominator, has the same
three local squareclasses as $r$, and has valuations at least
$4,4,2$ respectively. Also
\[
 0<R<\frac{102^2}{17^4}\frac{\ell^{1-2k}}{r^3}<1.
\]
The first inequality uses $(1+y^2)^2>y^4$; the last uses
$102^2<17^4$, $r\ge1$ and $k\ge1$.

Fix this integer input before choosing the elliptic indices. The set
of prime divisors of the numerator or denominator of $R$ is finite;
it includes $\ell$. Choose $H_0$ in a small compact neighborhood
of the real identity with $z(H_0)$ nonzero, and sufficiently deep in
the formal groups at this finite set. Choose a positive $s$, divisible
by the prescribed fixed modulus and by six, still deeper at these
places, and with $[s]P$ sufficiently close to the real identity. The
density of each prescribed compatible progression in the identity
component of the compact real closure supplies these choices. The
six normalized leaves can then be made arbitrarily small at the real
places and arbitrarily deep at the relevant finite places. Consequently
$W$ is a unit at every prime in the support of $R$ and is sufficiently
close to $1$ in the real absolute value. In particular one may arrange
$|W-1|<1/2$, so $0<RW^2<9/4<17$.

The scalar pole input $RW^2$ has square reduced denominator: outside
the fixed support its valuation is twice that of $W$, and on the
support $W$ is a unit. At $2,3,17$ it retains the required local
squareclasses and depths. The zero-preparation pole theorem therefore
supplies its rational certificate. Its first conic is covered by the
two specified rational charts, including its exceptional point.

The remaining simultaneous construction freezes these finite local
requirements in a progression before applying Sun's theorem. That
theorem allows avoidance at every fixed odd prime of this progression,
including $\ell$; avoidance at seven follows from its prescribed
Gaussian congruence. The reduction-order Chinese remainder conditions
remain compatible with the fixed progression. They also preserve the
chosen real neighborhood. Thus all three comparison poles and the
scalar pole are realized together, exactly with the assigned orders
$10,39$ and scalar order $20$. Larger fixed local depths require no new
rational coordinate. This completes the realization argument.
\end{proof}

\subsection{The exact polynomial and its bounds}

The source has $28$ retained rational coordinates: the input, $18$
point-abscissa coordinates, one slope, six comparison-pole coordinates,
and two scalar-pole coordinates. Its remaining finite choices are nine
monic quartic point roots, four monic quadratic pole roots, and four
bits choosing rational conic charts. The finite-algebra construction
eliminates the thirteen roots with one further rational coordinate.
Every zero and nonzero quadratic-root case is included. The maximum
finite-algebra rank is $4^9 2^4=2^{22}$.

\begin{theorem}
\label{thm:fixed-prime-29}
Under the four cited arithmetic inputs there is an explicitly specified
polynomial $F\in\mathbb Z[X,Y_1,\ldots,Y_{28}]$ such that
\[
 X\in\mathbb Z\quad\Longleftrightarrow\quad
 \exists (Y_1,\ldots,Y_{28})\in\mathbb Q^{28}\quad F(X,\boldsymbol Y)=0.
\]
It satisfies
\[
                  \deg F\le2^{174},\qquad
                  \max|\operatorname{coeff}(F)|\le2^{\,2^{177}}.
\]
The specification uses determinant, trace, resultant, coefficient
extraction and explicitly bounded finite products of integer polynomial
expressions. It has no finite product of elliptic-index exclusions.
\end{theorem}
\begin{proof}
The preceding local and height arguments prove soundness; the preceding
realization lemma and the simultaneous comparison construction prove
completeness. The guarded finite-algebra theorem applies to the stated
triangular monic relations and finite conic choices, leaving precisely
$28+1=29$ external coordinates. Its determinant and trace formulas have
integer coefficients and introduce no quantified matrix entries.

The coefficient directed acyclic graph, all monic relations, every
denominator guard, and the finite-product formula are supplied in
\texttt{polynomial-specification.json}. The independent checker
reconstructs every affine conversion from the rational source, verifies
all repeated-root branches, and recomputes the degree and coefficient
majorants by integer arithmetic. For auditability, its intermediate
bounds are
\[
 \begin{aligned}
 \text{matrix-entry degree}&\le4790038912,\\
 \log_2\text{(weighted operator norm)}&\le49766687378,\\
 \text{number of final factors}&=12292021203320302469120.
 \end{aligned}
\]
The resulting exact degree majorant has binary length $174$, and the
exact majorant for the logarithm to base two of the coefficient norm
has binary length $177$. The weighted companion-norm lemma and the elementary product, determinant and
resultant estimates in the finite-algebra theorem therefore give the
stated rounded bounds. This computation certifies the specified
algebraic construction and bounds; the four arithmetic theorems remain
its explicitly cited inputs.
\end{proof}
\subsection{Simultaneous completeness}

Fix $X\in\mathbb Z$. The input-map lemma gives an odd integer
$\beta>2401$. Lemma~\ref{lem:fixed-prime-completeness} proves that
the scalar contact $R_\beta$ is positive and less than one, has square
reduced denominator, and is square at $2,3,17$ with depths at least
$4,4,2$. Its finite support includes the fixed prime $\ell$.

Choose a real compact neighborhood for $H_0$ near, but away from, the
identity in both embeddings, with $f(H_0)$ small and $z(H_0)$ bounded
away from zero.  The positive orbit of every suitable finite-index
subgroup of $\mathbb ZP$ is dense in the identity component of its
real closure.  Thus one may choose a positive $h_*$ divisible by
$64$, and deep in every formal group over the finite support of
$R_\beta$, whose point lies in this neighborhood.  Choose a positive
integer $s$ divisible by six, still deeper at these places, with the fixed prime depth required in Lemma~\ref{lem:fixed-prime-index}, and with $[s]P$ sufficiently close to the real identity.
Set $a=s+1$, so $a$ is odd and $a\equiv1\pmod3$.
Take $S=[a]P$, and use $C_S=[s]P$ in the scalar certificate.
The normalized leaves of its norm tree are then small in the real
absolute value and sufficiently deep locally, uniformly for $H_0$ in
a smaller neighborhood.  Its root is a local unit and
$|W-1|<1/2$.

Odd Lucas labels are available explicitly.  Let
$U_0=0$, $U_1=1$, $U_{j+1}=aU_j-U_{j-1}$, and put
\[
              j=2(a-2),\qquad b=U_j/(a-2),\qquad t=U_{j+1}.
\]
Modulo $a-2$ one has $U_j\equiv j$, so $b$ is integral.
Modulo two the sequence has period $0,1,1$.  Since $a\equiv1\pmod3$,
$j\equiv1\pmod3$, so $b,t$ are odd.  They are positive and greater
than one.  The identity
$U_{j+1}^2-aU_jU_{j+1}+U_j^2=1$ gives $\Delta=0$.
The four root labels $\beta,a,t,b$ are therefore all odd and at least
three, and Lemma~\ref{lem:u29-chart} covers their points.

Temporarily set $H_0=[h]P$ and $U_i=[hn_i]P$.  Local continuity
provides a progression $h\equiv h_*\pmod L$, with $64\mid L$,
preserving the scalar units and fixed lower bounds for all comparison
tests at $2,3,17$.  Include the order of the finite component group
of the real closure in $L$.  Remove the finitely many comparison
pole loci from the real neighborhood and pass to a compact subset;
the tests are then uniformly bounded there.  These choices precede
the prime-pattern application.

Choose $D_0$ with $D_0^2\theta/A_r\in\mathbb Z$ for every contact
and scale, put $v=D_0^2w^2$, and choose $\lambda=u/v$.
The contact numerators $u-(\theta/A_r)v$ are a finite set of distinct
prime forms after repetitions are removed.  Apply the simultaneous
prime-pattern hypothesis beyond all fixed bad primes, with enough
divisibility of $w$ and a sufficiently large lower bound on $u/w^2$.
The curve has full rational two-torsion at every selected residue
field.  Because $P$ is a double, $d_\pm(p)/4$ annihilates its
reduction.  Define
\[
 U=\operatorname{lcm}_p(d_-(p)/4),\qquad
 V=\operatorname{lcm}_p(d_+(p)/4).
\]
The order and avoidance hypotheses give
$\gcd(U,V)=1$, $\gcd(U,L)=16$, and $\gcd(V,L)=1$.
Thus
\[
             h\equiv0\pmod U,\qquad h\equiv1\pmod V,
             \qquad h\equiv h_*\pmod L
\]
are compatible.  Their progression still has the required real
closure, so $H_0$ can lie in the chosen compact set.  At one place over
each contact prime the points $H_0,U_i$ are formal; at the conjugate
place $V_0,V_i$ are formal.  Both exact logarithmic ratios are $n_i$.
All sixteen tests therefore have their assigned orders, including the
direct Lucas polynomial since $\Delta=0$.

For completeness of each comparison pole, note the exact identity
\[
 R_r=\frac{\prod_{\theta\in\Theta}(u-(\theta/A_r)v)}
                     {v^5(A_r\lambda-1)^6}.
\]
The denominator is a rational square.  At the five contact primes
$v_p(R_r)=1$.  At either place over such a prime each inverse norm
factor has residue zero or one, and at the successful place every
residue is one.  Its reciprocal trace is thus a unit with residue
one or two.  Consequently $b_r$ has positive odd valuation only at
these primes and has square reduced denominator.  Their prescribed
congruences make $b_r$ square at $2,3,17$.

Uniformly on the compact real set,
$R_r\sim A_r^{-6}\lambda^{-1}$,
$W_r=1+O(\lambda^{-3})$, and
$b_r\sim(4A_r^6\lambda)^{-1}$.
Hence $0<b_r<17$ for a sufficiently large prescribed lower bound on
$\lambda$.  At $p\in\{2,3,17\}$ put $\nu=v_p(v)$ and
$\eta=v_p(A_r)$.  Sufficient divisibility of $w$ gives
$v_p(R_r)=\nu-6\eta$.  The nonconstant terms of the norm factors
have valuations at least $3\nu+(10m_j-8)\eta-2c_p$, where $c_p$
is the previously fixed lower-bound constant.  Their inverse trace
therefore has valuation $v_p(2)$, and
$v_p(b_r)=\nu-6\eta-2v_p(2)$.  These valuations can exceed every
required preparation threshold.  The pole hypothesis supplies the
three comparison certificates.

The scalar pole input retains square denominator and the specified
local squareclasses because $W$ is a unit at every prime in its fixed
support.  Also $0<R_\beta W^2<4<17$.  The same pole hypothesis
supplies its certificate with preparation zero.  All guards are met,
and the conic parametrization covers all these witnesses.  This proves
completeness for every integer input.

\subsection{The finite polynomial and its bounds}

The thirteen formal roots consist of nine point quartics and four
pole quadratics.  A rational equation $q^2=N/D$ with $D\ne0$ is
made monic by $Z=Dq$, so $Z^2-ND=0$.  Its zero-radicand branch is
$Z=0$ with $ND=0$; its other branch has $Z^2=ND$ and $Z\ne0$.
Together with the point derivative guards this makes every admitted
geometric point simple.  The rank is at most
\[
                              4^9\,2^4=2^{22}.
\]
Theorem~\ref{thm:explicit-guarded-finite-algebra} eliminates all thirteen
roots using one new rational variable $U$.  For completeness, its
literal recipe is as follows.  On each quadratic support branch,
put $m=13$ and let $n=4^9 2^k$, where $k$ is the number of nonzero
quadratic branches.  Let $G$ be the product of the source guards and $E_i$ the closed
zero-radicand conditions.  For every
$0\le\ell\le mn(n+1)/2$ and $1\le d\le n$, form
\[
 \lambda_\ell=1+\sum_{i=1}^m\ell^iz_i,\quad
 H(V)=\det(I-VM_{G\lambda_\ell})=\sum a_jV^j,
 \quad R_i(V)=\operatorname{Tr}(M_{E_i}\operatorname{adj}(I-VM_{G\lambda_\ell})).
\]
Put $H_d=\sum_{j=0}^d a_jV^j$,
$g=a_d\operatorname{Res}(H_d,H'_d)$, $L=a_d^2g$, and take the factor
\[
 \left(1+\sum_{j=1}^d a_{d-j}a_d^{2j-1}g^jU^j\right)^2
       +\sum_{j>d}a_j^2
       +\sum_i\left((LU)^{n-1}R_i(1/(LU))\right)^2.
\]
Coefficient reversal makes the last expressions polynomials at
$LU=0$.  Multiply these factors over the displayed ranges, the
$2^4$ quadratic masks, and all $2^4$ original finite choices.
This is a completely specified integer polynomial in the retained
coordinates and $U$, with no formal root left as a free variable.

\subsection{Weighted companion norms}
\label{subsec:weighted-companion-norms}

The coefficient estimate can be sharpened without changing the
polynomial. The useful norm on a monic finite algebra assigns different
positive scales to its basis monomials. We describe this choice and its
compatibility with all determinant and trace operations.

Let $B=\mathbb R[X_1,\ldots,X_r]$, equipped with the coefficient norm
$\|f\|_1=\sum_\alpha|[X^\alpha]f|$. This norm is submultiplicative.
Consider the free iterated monic algebra
\[
 A_0=B,\qquad
 A_i=A_{i-1}[z_i]\big/
 \left(z_i^{d_i}+\sum_{j<d_i}c_{ij}z_i^j\right).
\]
Fix the usual basis of monomials $z_1^{e_1}\cdots z_i^{e_i}$,
$0\le e_j<d_j$. Coefficients $c_{ij}$ belong to $A_{i-1}$.
Integer coefficients are a special case of this real coefficient ring;
no assertion about zeros over a larger field is being made.

\begin{lemma}\label{lem:weighted-companion-norm}
Suppose the coefficient multiplication operators in $A_{i-1}$ have
bounds $\|M_{c_{ij}}\|\le2^{b_{ij}}$, with $b_{ij}\ge0$.
Let $q_i$ be the number of nonzero coefficients, or any upper bound
for that number. Take $q_i\ge1$, and put
\[
 \gamma_i=
 \max\left(\{0\}\cup
   \left\{\left\lceil\frac{b_{ij}}{d_i-j}\right\rceil:
                  c_{ij}\ne0\right\}\right)
       +\lceil\log_2 q_i\rceil.
\]
On $A_i$ use the norm
\[
 \left\|\sum_{j=0}^{d_i-1}a_jz_i^j\right\|
               =\sum_{j=0}^{d_i-1}2^{j\gamma_i}\|a_j\|.
\]
Then $\|M_{z_i}\|\le2^{\gamma_i}$. Multiplication by an element
of $A_{i-1}$ retains its preceding operator-norm bound.
\end{lemma}
\begin{proof}
For a basis block $z_i^j$ with $j<d_i-1$, multiplication by $z_i$
shifts to the next block and multiplies its norm weight by
$2^{\gamma_i}$. On the last block, the monic relation gives a column
of norm at most
\[
 \sum_{j<d_i}\|M_{c_{ij}}\|2^{(j-d_i+1)\gamma_i}
 =2^{\gamma_i}
       \sum_{j<d_i}\|M_{c_{ij}}\|2^{-(d_i-j)\gamma_i}.
\]
Write $u=\lceil\log_2 q_i\rceil$. Every nonzero summand in the last
sum is at most $2^{-(d_i-j)u}\le2^{-u}$, and there are at most
$q_i$ summands. The sum is therefore at most one. These are all
block columns of the companion matrix, proving the desired bound.
An earlier multiplication operator acts independently and identically
on every block. Its block weights cancel, so its preceding norm bound
is unchanged.
\end{proof}

Starting with $\|f\|_1$ on $B$, the lemma defines a norm throughout
the monic tower. The triangle and product inequalities bound any
coefficient expression recursively. In particular, every integer
constant $a$ has norm $|a|$, every base coordinate has norm one, and
an arithmetic expression with recorded logarithmic bound $b$ has
operator norm at most $2^b$.

For completeness, these weighted operator bounds can be used in the
coefficient estimates for the characteristic and trace polynomials
without introducing any extra factor for the size of the basis scales.
Let $D$ be the constant diagonal matrix whose diagonal entry for a
basis monomial is
\[
                         \prod_i2^{e_i\gamma_i}.
\]
The weighted maximum column norm of $M$ is exactly the ordinary maximum
column sum of the coefficient norms of $DMD^{-1}$. Simultaneous
conjugation leaves
\[
 \det(I-VM),\qquad
 \operatorname{Tr}\bigl(M_E\operatorname{adj}(I-VM)\bigr)
\]
unchanged. The adjugate identity follows first for invertible matrices
from $\operatorname{adj}(A)=\det(A)A^{-1}$ and then for all matrices
by polynomial identity. Thus every determinant, trace and resultant
estimate already stated may be applied to the conjugated matrices.
The coefficients of the resulting scalar polynomials are exactly the
original coefficients. In particular, no additional basis-height
term is needed. The weighted degree calculation is separate and is
unchanged.

If a quadratic generator is set to zero in one of the finite branches,
its norm is zero. The recursive estimates for all later expressions
remain valid after that substitution. Thus a single table also bounds
all of the zero and nonzero quadratic cases.

For the construction with $29$ variables, this gives the generator
logarithmic norm bounds
\[
 \underbrace{5,\ldots,5}_{9\text{ entries}},\quad
 89263067,\quad1539436917,\quad2988679217,\quad19861198427.
\]
The weighted degree table is unchanged. The multiplication operator
used in the determinant has logarithmic norm bound $49766687378$.
The exact final bound for the logarithm to base two of the coefficient
norm is
\[
 180551808614174879228877281249630444945389422609694720
                                                        <2^{177}.
\]
Consequently the same polynomial satisfies
\[
                    \max|\operatorname{coeff}(F)|\le2^{\,2^{177}}.
\]
The independent checker recomputes these bounds from the full integer
coefficient graph. A separate equality certificate verifies that
removing the bound annotations leaves precisely the same mathematical
specification. The degree bound remains $2^{174}$.
\subsection{Explicit degree and coefficient bounds for the specification}
\label{subsec:finite-algebra-bounds}

The following bounds concern the expanded polynomial defined by the
indexed determinant formula. They do not require expanding its matrices
or its final finite product. For a polynomial $f$, let $\|f\|_1$ denote
the sum of the absolute values of its coefficients. Then
$\|fg\|_1\le\|f\|_1\|g\|_1$ and
$\|f+g\|_1\le\|f\|_1+\|g\|_1$.

Give every base coordinate weight one, and assign a positive integer
weight $w_i$ to each algebra generator $z_i$. Choose the weights so that
every nonleading term of its monic relation has weight at most
$e_iw_i$. Reduction by these relations never increases weighted degree.
Writing $W=\sum_i(e_i-1)w_i$, the entries of a multiplication matrix
$M_f$ therefore have base degree at most $\operatorname{wdeg}(f)+W$:
the input basis monomial has weight at most $W$, and the output basis
monomial has nonnegative weight.

Use the weighted multiplication-operator norms of
Lemma~\ref{lem:weighted-companion-norm}. Their triangle and product
inequalities give the coefficient bounds $\|M_{z_i}\|\le2^{b_i}$
by recursion. The weights used for polynomial degrees are separate:
the point-quartic generators have degree weight two, while each
quadratic generator may be assigned half its constant coefficient's
weighted degree, rounded upward. Its zero branch satisfies the same
bounds. A single table therefore covers all $2^4$ quadratic cases.
The constant diagonal conjugation proved in the preceding subsection
justifies using these smaller operator norms in every scalar
determinant and trace estimate below.

Put $n=2^{22}$ and $m=13$, and
$T=m n(n+1)/2$ for the largest separating parameter. Suppose the
recorded bounds give
\[
 \operatorname{wdeg}(G)\le w_G,\quad \|M_G\|\le2^{b_G},\qquad
 \operatorname{wdeg}(E_i)\le w_E,\quad \|M_{E_i}\|\le2^{b_E}.
\]
Set
\begin{align*}
 b_\lambda&=\lceil\log_2(m+1)\rceil
       +\max_i\{i\lceil\log_2 T\rceil+b_i\},\\
 \alpha&=w_G+\max_i w_i+W,\qquad \beta=b_G+b_\lambda,\\
 \alpha_E&=w_E+W.
\end{align*}
These bound the entry degrees and operator norm of
$M_{G\lambda_\ell}$, uniformly for all prescribed $\ell$ and all
source branches.

The coefficients of $H(V)=\det(I-VM)$ have base degree at most
$A_H=n\alpha$, and its full coefficient norm is at most
$2^{B_H}$ with $B_H=n(\beta+1)$. Indeed the coefficient norm of a
determinant is bounded by the product of its column sums. This bound
already includes the sum over determinant permutations; no additional
factorial is needed. By multilinearity in columns,
\[
 R_i(V)=\left.\frac{\partial}{\partial t}
       \det(I-VM+tM_{E_i})\right|_{t=0}
\]
has base coefficient degree at most
$A_R=(n-1)\alpha+\alpha_E$ and coefficient norm at most $2^{B_R}$,
where
\[
 B_R=\lceil\log_2 n\rceil+b_E+(n-1)(\beta+1).
\]
This accounts for the $n$ possible columns containing $M_{E_i}$.

For each degree case $d\le n$, the Sylvester matrix for
$\operatorname{Res}(H_d,H'_d)$ has $2d-1$ rows. Every row has
coefficient norm at most $d\|H\|_1$. Hence uniform bounds for the
reciprocal scale $L=a_d^3\operatorname{Res}(H_d,H'_d)$ are
\begin{align*}
 A_L&=(2n+2)A_H,\\
 B_L&=(2n+2)B_H+(2n-1)\lceil\log_2 n\rceil.
\end{align*}
The reciprocal residual has degree at most $n(A_L+1)$ and coefficient
norm at most $2^{nB_L+\lceil\log_2(n+1)\rceil}$. The reversed trace
residuals have degree at most $A_R+(n-1)(A_L+1)$ and coefficient norm
at most $2^{B_R+(n-1)B_L}$. Together with the high-coefficient
conditions $a_k=0$, these give a bound for the sum of squares by
doubling the largest degree, doubling the largest logarithmic norm
bound, and adding $\lceil\log_2(n+q+1)\rceil$ to the latter.
For the reciprocal estimate, no division of coefficient norms is used:
its $j$th coefficient is the polynomial
$a_{d-j}a_d^{3j-1}\operatorname{Res}(H_d,H'_d)^j$.
Its degree is at most $jA_L$ and its logarithmic norm bound is at most
$j(3B_H+B_{\rm Res})=jB_L$.

The number of final factors is also explicit. Put $n_0=4^9$ and $q=4$.  For a mask with $k$ nonzero quadratic branches, the rank is
$n_k=n_0 2^k$.  There are $\binom qk$ such masks, each with
$n_k(1+mn_k(n_k+1)/2)$ separating-parameter and degree choices.
Including the $2^4$ original finite parameter choices gives exactly
\[
 B_{\mathrm{tot}}=2^4\left[
 n_0 3^q+\frac m2\bigl(n_0^3 9^q+n_0^2 5^q\bigr)\right]
\]
factors. Its exact value is
\[
                     B_{\mathrm{tot}}=12292021203320302469120.
\]
Multiplying the uniform branch-degree and logarithmic coefficient-norm
bounds by this integer bounds the final product. All the required
weights and norms are computed by finite recursions on the supplied
coefficient graph. The exact integer majorants give
\[
                         \deg F\le2^{174},\qquad
                         H(F)\le2^{\,2^{177}}.
\]
These estimates require no expanded matrix or monomial table. The
independent checker evaluates the recursions and compares every
integer bound in the saved specification.
\section{The input map and its local squareclasses}
\label{sec:scalar-correction-free}

Put
\[
                         \beta(X)=216X^4+450459
\]
and retain the five-contact polynomial
\[
 r(T)=T(T-1)(T-2401)(32T+343)(2T-1).
\]
Define
\[
                         r_X=r(\beta(X)).
\]


\begin{lemma}\label{sn:integral-input}
For $X\in\mathbb Q$ one has
\[
                 \beta(X)\in\mathbb Z
                 \quad\Longleftrightarrow\quad X\in\mathbb Z.
\]
For every real $X$, $\beta(X)\ge450459>2401$, so $r_X>0$
.  If $p>76176$ and $v_p(\beta(X))<0$,
then
\[
 v_p(\beta(X))=4v_p(X),\qquad
 v_p(r_X)=5v_p(\beta(X)).
\]
In particular, both valuations are even.
\end{lemma}
\begin{proof}
If $X$ is an integer, so is $\beta(X)$.  Conversely, write
$X=a/b$ in lowest terms with $b>0$.  Integrality of $\beta(X)$
gives $b^4\mid216=2^3\cdot3^3$.  Neither prime exponent
admits a nontrivial fourth power, so $b=1$.

The real inequalities follow from the fourth power and the
five positive factors of $r(T)$ when $T>2401$.  At a prime
in the stated range, a negative valuation of $\beta$ forces
$v_p(X)<0$.  The unique dominant terms in $\beta$ and
$r(\beta)$ are respectively $216X^4$ and $64\beta^5$.
These identities prove the claimed valuation formulas.
\end{proof}

\begin{lemma}\label{sn:local-squares}
If $X\in\mathbb Z$, then $r_X$ is a nonzero square in
$\mathbb Q_2$, $\mathbb Q_3$, and $\mathbb Q_{17}$.
More precisely,
\[
 v_2(r_X)=2,\quad r_X/4\equiv1\pmod8,
 \qquad
 v_3(r_X)=2,\quad r_X/9\equiv1\pmod3,
\]
and at seventeen either $v_{17}(r_X)=0$ with square-unit
residue, or
\[
                    v_{17}(r_X)=2,\qquad
                    r_X/17^2\equiv9\pmod{17}.
\]
\end{lemma}
\begin{proof}
Since $216$ is divisible by eight and by twenty-seven,
\[
                    \beta\equiv3\pmod8,
                    \qquad\beta\equiv18\pmod{27}.
\]
For the two even factors of $r(\beta)$ one has
\[
 v_2(\beta-1)=v_2(\beta-2401)=1,
 \qquad
 \frac{\beta-1}{2}-\frac{\beta-2401}{2}=1200\equiv0\pmod8.
\]
Their odd parts therefore have square product congruent to one
modulo eight.  The remaining odd factors have product
$3\cdot7\cdot5\equiv1\pmod8$.  This proves the assertions
at two.  At three, only the factor $\beta$ is divisible by
three; it has valuation two.  Reduction after division by nine
gives
\[
 \frac{r(\beta)}9\equiv
       2\cdot2\cdot2\cdot1\cdot2\equiv1\pmod3.
\]

At seventeen the fourth powers are exactly $0,1,4,13,16$.
Since $216\equiv12$ and $450459\equiv10\pmod{17}$, all
possibilities are given by the table
\[
\begin{array}{c|rrrrr}
 X^4&0&1&4&13&16\\\hline
 \beta&10&5&7&13&15\\
 r(\beta)&0&15&2&13&2.
\end{array}
\]
The nonzero residues are squares: $13=8^2$, $15=7^2$, and
$2=6^2$ modulo seventeen.

In the remaining case $17\mid X$.  The exact identity
\[
 32\cdot450459+343=17^2\cdot49879,
 \qquad 49879\equiv1\pmod{17},
\]
and $v_{17}(X^4)\ge4$ show that the factor
$32\beta+343=6912X^4+17^2\cdot49879$ has valuation exactly
two and unit residue one.  The other four factors are units,
with product
$10\cdot9\cdot6\cdot2\equiv9\pmod{17}$.
This proves the last assertion.  The usual local square criteria
at two and at odd primes now prove all three square statements.
\end{proof}

\section{Number-field scope}

All displayed algebraic identities remain valid in characteristic
zero.  Positivity of the contact and norm root also holds in a
number field with a real embedding, after applying that embedding.
The local squareclass table, however, concerns rational integer
inputs and the three specified completions of $\mathbb Q$.
It does not establish the corresponding preparation assertions
for every algebraic integer in a larger field.

Even the integrality equivalence of Lemma~\ref{sn:integral-input}
does not automatically become an equivalence for rings of
integers.  Let $x_0$ be a real fourth root of $1/108$ and
$L=\mathbb Q(x_0)$.  Then
\[
                       \beta(x_0)=450461\in\mathbb Z.
\]
But $x_0$ is not an algebraic integer, since $x_0^4=1/108$.
Thus this elementary input map
alone does not transfer a definition to $\mathcal O_L$.
No assertion that the complete rational defining polynomial
defines every $\mathcal O_L$ is made or follows from this
normalization.


\section{Exact specification and reproducibility}
The directory \texttt{computation/few\_variables/forced\_prime29/}
contains the rational source, its producer, the integer polynomial
specification, and independent checkers. The reproduction script
rebuilds the source and specification byte for byte, checks all monic
relations and denominator guards, and recomputes the degree and
coefficient estimates. Separate integer checks certify the chosen
prime, all five fourth-power exclusions, and the strict index margin.

The determinant and resultant matrices and the final finite product
are specified exactly; their full expansions are not stored. The
global arithmetic inputs remain the four stated premises of the
projection theorem. No computational certificate is used as a
replacement for one of those theorems.
\begin{thebibliography}{99}
\small
\bibitem{RepositoryV9}
E.~Glazer and contributors, \emph{Diophantine definition of the integers
in the rationals}, source repository, coefficient baseline
\texttt{14e8f601fe0d9173d4a5f084aedc3e07066a7688}, 7 November 2026.
The odd-product identity, direct targets and asymmetric-depth
construction are in \texttt{docs/quartic-523-540/addition.tex}; the
repository attributes the odd-product contribution to G.~Conant,
developed with ChatGPT.
\url{https://github.com/elliotglazer/diophantine-z-q/tree/14e8f601fe0d9173d4a5f084aedc3e07066a7688}.


\bibitem{RepositoryV11}
E.~Glazer and contributors, \emph{Diophantine definition of the integers
in the rationals}, corrected arithmetic proof, Version~11,
commit \texttt{650b298a2a18fa956f8859f3ca80e9d7c8cedd82}.
\url{https://github.com/elliotglazer/diophantine-z-q/tree/650b298a2a18fa956f8859f3ca80e9d7c8cedd82/docs/manuscript-revisions/version11}.

\bibitem{Sun2026Rational86}
Z.-W.~Sun, \emph{86 unknowns for Hilbert's tenth problem over
the rationals}, arXiv:2610.20691v2, 19 October 2026,
especially Theorem 4.1 and Sections 10.3--10.6.
\url{https://arxiv.org/abs/2610.20691v2}.


\bibitem{Silverman2009}
J.~H.~Silverman, \emph{The Arithmetic of Elliptic Curves},
second edition, Graduate Texts in Mathematics 106, Springer, 2009.
\url{https://doi.org/10.1007/978-0-387-09494-6}.


\end{thebibliography}
\end{document}
